\section{Performance Evaluation}
\label{sec:evaluation}

This section evaluates the proposed cooperative assignment-and-routing
framework in an electric waste-collection digital twin constructed from a
real urban road network. The experiments address four questions: (i) whether
the proposed method consistently reduces operating energy under the hourly
traffic conditions of the service area, (ii) how the exact number of active
vehicles changes energy consumption and completion time, (iii) whether
replanning is beneficial when localized traffic disruptions evolve during
operation, and (iv) whether the solution quality is retained as the number of
service bins increases.

\begin{table}[t]
    \centering
    \caption{Simulation Parameters}
    \small
    \setlength{\tabcolsep}{2pt}
    \renewcommand{\arraystretch}{1.1}
    \begin{tabular}{lr}
    \toprule
    \textbf{Parameter} & \textbf{Value / Description} \\
    \midrule
    Number of smart bins, $N$ & $84$ \\
    Scalability values of $N$ & $84,100,\ldots,500$ \\
    Baseline road network & OpenStreetMap/SUMO \\
    Time-dependent traffic & Hourly SUMO/TOPIS travel times \\
    Number of vehicles, $K$ & $10$, $\{6,\ldots,14\}$ for sweep \\
    Scalability fleet size & approximately $N/8.4$ \\
    Curb mass, $m_0$ & 7,000~kg \\
    Payload capacity, $Q$ & 2,000~kg \\
    Maximum battery capacity, $B_{\max}$ & 180~kWh \\
    Minimum battery reserve, $B_{\min}$ & 15~kWh \\
    Node weight, $w_i$ & $\mathcal{U}(50,250)$~kg \\
    Rolling coefficient, $C_{\mathrm{rr}}$ & 0.010 \\
    Drag coefficient, $C_{\mathrm{d}}$ & 0.60 \\
    Frontal area, $A_{\mathrm{f}}$ & 5.5~$\mathrm{m}^2$ \\
    Air density, $\rho$ & 1.225~$\mathrm{kg/m^3}$ \\
    Propulsion efficiency, $\eta_{\mathrm{drv}}$ & 0.90 \\
    Regenerative efficiency, $\eta_{\mathrm{reg}}$ & 0.60 \\
    En-route auxiliary power, $P_{\mathrm{aux}}$ & 3~kW \\
    Stop-and-go threshold, $\Delta T$ & $0.25T^{\mathrm{ff}}$ \\
    Stop-and-go coefficient, $\kappa_{\mathrm{sg}}$ &
        $4.93\times10^{-6}$~kWh/kg \\
    Service time per bin, $T_s$ & 120~s/bin \\
    Stationary auxiliary power, $P_{\mathrm{svc}}$ & 2~kW \\
    Lifting/compaction energy, $e_i^{\mathrm{lc}}$ & 0.12~kWh/bin \\
    Payload state interval, $\Delta\ell$ & 100~kg \\
    \bottomrule
    \end{tabular}
    \label{tab:parameters}
\end{table}

\subsection{Simulation setup}

The simulation parameters are summarized in Table~\ref{tab:parameters}. The
evaluation uses the Seongbuk-gu road network in Seoul. The network is
generated from OpenStreetMap, represented as a directed SUMO graph, and
augmented with road grade and hourly link travel time. The base instance
contains $N=84$ measured waste-bin locations. For every transition, all
methods use the same load-, grade-, speed-, and time-dependent energy model in
Section~II, including auxiliary, stop-and-go, service, and
lifting/compaction energy. Therefore, the reported energy is battery energy
rather than distance or a surrogate routing cost.

Unless otherwise stated, the base experiments use $K=10$ vehicles and demand
seed 7. Every returned plan must service every bin exactly once, respect the
2,000~kg vehicle capacity and 15~kWh battery reserve, return every vehicle to
the depot, and assign at least one bin to each of the exactly $K$ active
vehicles. Infeasible solutions are not assigned a finite objective value.

The proposed method is compared with a state-aware nearest-neighbor method
(NN), ant-colony optimization (ACO), particle-swarm optimization (PSO), and a
genetic algorithm (GA). ACO uses 60 ants and 150 iterations, PSO uses 60
particles and 200 iterations, and GA uses a population of 100 and 500
generations. Their candidate permutations are decoded into exactly $K$
nonempty, capacity- and battery-feasible routes by the same dynamic-programming
split procedure. Thus, the comparison changes the search method while keeping
the physical evaluator and feasibility rules fixed. The MILP implementation
is retained for small-instance correctness tests, but is not used as an
operational-size baseline. If the loopy message updates do not satisfy the
strict residual test within the prescribed round budget, the best feasible
incumbent is reported; consequently, the experiments do not claim global
optimality.

\subsection{Hourly traffic and route comparison}

\begin{figure}[!t]
    \centering
    \subfloat[Proposed: 164.679~kWh]{
        \includegraphics[width=0.78\columnwidth]{figures/exp1_kmedoids}
    }

    \vspace{0.2cm}
    \subfloat[ACO: 172.225~kWh]{
        \includegraphics[width=0.78\columnwidth]{figures/exp1_ap}
    }

    \vspace{0.2cm}
    \subfloat[GA: 188.522~kWh]{
        \includegraphics[width=0.78\columnwidth]{figures/exp1_proposed}
    }
    \caption{Fleet-route comparison on the same Seongbuk-gu instance at
    11:00. All methods service the same 84 bins with exactly 10 nonempty and
    feasible vehicle routes. Colors distinguish vehicle routes, black circles
    indicate smart bins, and the star indicates the depot.}
    \label{fig:route-comparison}
\end{figure}

The first experiment fixes the bin demands and evaluates all methods for the
24 traffic snapshots from 00:00 to 23:00. Figure~\ref{fig:route-comparison}
shows the routes obtained at 11:00 by the proposed method and the two strongest
stochastic baselines in this snapshot. The proposed method uses 7.547~kWh less
than ACO and 23.843~kWh less than GA. Its 3,261.6~s makespan is nevertheless
longer than the 2,809.1~s and 2,806.8~s obtained by ACO and GA, respectively.
This example illustrates that the optimized objective is fleet energy, not
makespan, and that the two metrics need not rank route sets identically.

\begin{figure}[!t]
    \centering
    \subfloat[Fleet energy consumption]{
        \centering
        \includegraphics[width=0.95\columnwidth]{figures/hourly_energy_ieee.pdf}
    }

    \subfloat[Mean active-vehicle completion time]{
        \centering
        \includegraphics[width=0.95\columnwidth]{figures/hourly_vehicle_completion_time_ieee.pdf}
    }
    \caption{Performance over the 24 hourly traffic snapshots provided
    by~\cite{seoul2026}. The demand realization and fleet size are fixed across
    all hours and methods.}
    \label{fig:hour-sweep}
\end{figure}

As shown in Fig.~\ref{fig:hour-sweep}, the proposed method has the lowest
energy in all 24 matched hourly comparisons. Its mean energy is 116.893~kWh,
compared with 121.961~kWh for ACO, 130.857~kWh for GA, 141.951~kWh for PSO,
and 148.240~kWh for NN. This corresponds to reductions of 4.16\%, 10.67\%,
17.65\%, and 21.15\%, respectively, when the ratio of the 24-hour means is
used. Even the smallest hourly margin over ACO is 1.623~kWh.

Although makespan is not included in Fig.~\ref{fig:hour-sweep}, its hourly
variation provides a complementary measure of operational robustness. ACO's
makespan increases from 2,809.1~s at 11:00 to 4,126.4~s at 12:00, a 46.9\%
increase in one hour, while PSO increases from 2,818.4~s at 20:00 to
4,062.7~s at 21:00, corresponding to 44.2\%. GA also reaches 3,519.9~s at
14:00 after recording 2,977.5~s at 13:00. The proposed method does not exhibit
these 3,500--4,000~s excursions: its makespan remains between 2,500.4 and
3,261.6~s over all 24 snapshots. Its temporal standard deviation is also the
smallest, at 192.5~s, compared with 413.7~s for ACO, 326.9~s for PSO,
258.9~s for GA, and 200.9~s for NN. This reduced sensitivity of the longest
route indicates that the cooperative exchange between assignment and routing
also limits severe route imbalance under changing traffic, even though
makespan is not the optimized objective.

This distinction is important because minimizing fleet energy does not imply
minimum makespan at every hour; for example, GA has a slightly lower
24-hour mean makespan than the proposed method. Nevertheless, when completion
time is averaged over the active vehicles rather than determined only by the
single longest route, the proposed method requires 1,631.8~s per vehicle,
compared with 1,679.2~s for ACO, 1,743.5~s for GA, 1,817.2~s for PSO, and
1,864.0~s for NN. Thus, although the primary objective is energy consumption,
the resulting assignments and routes also reduce the average time required by
an individual vehicle. The proposed energy itself ranges from 78.198~kWh at
04:00 to 168.065~kWh at 12:00. Its nonmonotonic hourly pattern reflects the
spatial distribution of link-specific congestion and the stop-and-go term,
which cannot be explained by a single network-wide average speed.

\subsection{Online adaptation to localized disruptions}

The second experiment begins at 18:00 and uses ten paired traffic seeds. For
each seed, approximately 5\% of the bins are selected as fixed disruption
centers. Traffic remains normal for 10~min; speeds on edges within a 500~m
road-distance radius then decrease linearly from 100\% to 35\% over 10~min,
remain at 35\% for 10~min, and recover linearly over 25~min. Every method
observes the same trace and may replan every 5~min.

Three policies are distinguished. Frozen static retains both the initial bin
order and its precomputed edge paths. Live-navigation static retains the
initial bin order but follows minimum-energy paths under the evolving traffic.
Adaptive additionally reoptimizes the remaining bin assignments and orders.
Vehicles continue moving while planning is performed, so planning time is not
converted into stationary energy or operating delay; it is measured
separately. An adaptive candidate is accepted only if exact evaluation under
the currently observed traffic snapshot gives lower residual energy than the
incumbent. No traffic forecast is supplied.

\begin{figure}[!t]
    \centering
    \subfloat[Final operating energy]{
        \centering
        \includegraphics[width=0.95\columnwidth]{figures/online_final_energy_bar_ieee.pdf}
    }

    \subfloat[Snapshot-estimated total energy]{
        \centering
        \includegraphics[width=0.95\columnwidth]
        {figures/online_snapshot_total_energy_ieee.pdf}
    }

    \subfloat[Performance table\label{tab:disruption}]{
    \centering
    \footnotesize
    \setlength{\tabcolsep}{2.5pt}
    \begin{tabular}{lcccc}
        \toprule
        \multirow{2}{*}{\textbf{Method}} & \multicolumn{2}{c}{\textbf{Operational energy (kWh)}} & \textbf{Saving} & \textbf{Replanning}\\
        \cmidrule(lr){2-3}
        & \textbf{Frozen} & \textbf{Adaptive} & \textbf{(\%)} & \textbf{latency (min.)}\\
        \midrule
        GA & 190.50 ($\pm$14.43) & 184.70 ($\pm$13.69) & 3.0 & 44.07 ($\pm$3.65)\\
        PSO & 220.30 ($\pm$22.02) & 180.04 ($\pm$14.27) & 18.3 & 32.10 ($\pm$1.56)\\
        ACO & 171.97 ($\pm$8.76) & 159.20 ($\pm$7.24) & 7.4 & 10.91 ($\pm$0.60)\\
        NN & 238.05 ($\pm$25.04) & 196.91 ($\pm$18.35) & 17.3 & 0.89 ($\pm$0.15) \\
        Proposed & 168.87 ($\pm$14.33) & 158.31 ($\pm$9.92) & 6.3 & 9.76 ($\pm$0.73)\\
        \bottomrule
    \end{tabular}
    }
    \caption{Online coordination performance under localized road disruptions.}
    \label{fig:online_disruption}
\end{figure}

Figure~\ref{fig:online_disruption} and Table~\ref{tab:disruption} show that
the proposed adaptive method has the lowest mean final energy, 158.313~kWh.
It is 0.885~kWh below adaptive ACO and 21.725--38.597~kWh below adaptive PSO,
GA, and NN. Relative to its own frozen-static execution, proposed replanning
saves 10.555~kWh (6.04\%) and wins in all ten paired seeds. Its cumulative
post-departure replanning time is 9.76~min, lower than ACO, PSO, and GA but
higher than NN.

Adaptive planning is not uniformly better than live-navigation static. For
the proposed method, adaptive execution is 0.362~kWh higher on average and
wins in six of ten seeds; GA shows a larger degradation. This does not
contradict the acceptance test, which compares residual plans only under the
current snapshot. A route accepted during degradation may lose its apparent
advantage as traffic subsequently recovers. Thus, the experiment supports
adaptive assignment under unanticipated path disruption, but also identifies
the limitation of myopic replanning without adding a traffic-forecast model.

\subsection{Effect of the exact number of active vehicles}

\begin{figure}[!t]
    \centering
    \subfloat[Normalized energy gap]{
        \centering
        \includegraphics[width=0.95\columnwidth]{figures/k_sweep_normalized_energy_ieee.pdf}
    }

    \subfloat[Normalized makespan gap]{
        \centering
        \includegraphics[width=0.95\columnwidth]{figures/k_sweep_normalized_makespan_ieee.pdf}
    }
    \caption{Exact-$K$ sweep at 04:00, 11:00, and 19:00 with $N=84$ and
    11,700~kg total demand. Each line is the three-hour mean gap to the best
    method-and-$K$ value within the same hour, and each shaded region spans the
    minimum and maximum over the three hours.}
    \label{fig:k-sweep}
\end{figure}

The third experiment varies the exact number of active vehicles from 6 to 14.
The total demand is fixed at 11,700~kg, so six vehicles provide only 300~kg of
aggregate spare capacity, whereas larger fleets introduce additional
depot-to-service and return legs. For a metric $z$, Fig.~\ref{fig:k-sweep}
reports
\begin{align}
g_{m,K,h}=100\left(
\frac{z_{m,K,h}}{\min_{m',K'}z_{m',K',h}}-1\right),
\label{eq:k-sweep-normalization}
\end{align}
followed by the mean and range of $g_{m,K,h}$ over the three evaluated hours.

The proposed method consumes less energy than every baseline in all 27
matched hour-and-$K$ cases. Averaged over these cases, it reduces energy by
5.55\% relative to ACO, 10.62\% relative to GA, 18.65\% relative to PSO, and
22.30\% relative to NN. Its three-hour mean energy is minimized at $K=6$
(107.100~kWh); the individual-hour minima occur at $K=7$ for 04:00 and at
$K=6$ for 11:00 and 19:00. Energy then generally increases as more vehicles
repeat fixed depot-access and return overhead.

The completion-time result exhibits a different optimum. The proposed
method's lowest three-hour mean makespan is 2,940.0~s at $K=9$, rather than at
its energy-optimal $K=6$. Moreover, makespan is not monotone in $K$, because
adding a vehicle changes both the assignment and the longest route. The
experiment therefore identifies an operational tradeoff instead of imposing
a convex fleet-size curve: a smaller feasible fleet minimizes energy in this
instance, whereas a moderately larger fleet improves completion time.

\subsection{Scalability with the service-network size}

\begin{figure}[!t]
    \centering
    \subfloat[Energy per serviced bin]{
        \centering
        \includegraphics[width=0.95\columnwidth]{figures/scalability_energy_per_bin_ieee.pdf}
    }

    \subfloat[Makespan]{
        \centering
        \includegraphics[width=0.95\columnwidth]{figures/scalability_makespan_ieee.pdf}
    }
    \caption{Scalability with nested service-node sets and approximately 8.4
    bins per vehicle. Lines show seed means and shaded regions show one sample
    standard deviation.}
    \label{fig:scalability-test}
\end{figure}

The final experiment, summarized in Fig.~\ref{fig:scalability-test}, uses nested instances
$N\in\{84,100,120,\ldots,300,350,400,450,500\}$.
The fleet size is adjusted to keep approximately 8.4 serviced bins per
vehicle, so increasing $N$ enlarges the number of coupled assignment edges
while keeping the typical vehicle-level trellis size comparable. The demand
seed changes both the waste realization and the randomized stochastic search,
whereas the measured 84-bin locations remain fixed and the additional bins
are drawn once from a seeded nested set of feasible SUMO road locations.

For the ten-seed portion currently complete through $N=300$, every method
returns feasible solutions. The proposed method has the lowest mean energy per
bin at all 12 completed sizes. Its value changes from
$1.963\pm0.039$~kWh/bin at $N=84$ to $1.995\pm0.042$~kWh/bin at $N=300$ and
remains between 1.936 and 2.020~kWh/bin over the completed range. Averaged
over the 12 sizes, the corresponding reduction is 2.85\% relative to ACO,
10.73\% relative to GA, 14.80\% relative to PSO, and 16.80\% relative to NN.
The stable normalized energy indicates that total fleet energy grows mainly
with the amount of serviced work rather than through a degradation of
per-bin solution quality.

The makespan curve is less regular because it is determined by the single
longest vehicle route and is not the optimized objective. For the proposed
method, the mean makespan is 3,371.0~s at $N=84$ and 4,284.2~s at $N=300$,
with intermediate changes caused by the road geometry, demand realization,
and discrete fleet-size adjustment. Runtime is retained in the result files
for reproducibility, but is not used as a primary figure because it combines
hardware-dependent solver time with cost-oracle preparation; the analytical
growth of the assignment and trellis updates is discussed separately in the
complexity analysis.

Together with the hourly and fleet-size experiments, this scaling behavior
shows that the assignment--trellis cooperation improves energy efficiency
without sacrificing feasibility or producing the pronounced longest-route
delays observed in several baselines. The online results further show that the
same coordination can adapt assignments after an unanticipated disruption,
while the comparison with live navigation clarifies that this benefit depends
on the information available at the replanning epoch. The proposed framework
therefore combines lower fleet energy, stable vehicle-level workload, and
adaptive operation within the same message-passing structure rather than
treating these observations as separate design objectives.
